\section*{8. REFLECTION ON APPLYING SOLID}

\phantomsection
\addcontentsline{toc}{section}{8. REFLECTION ON APPLYING SOLID}

This chapter looks back on how the SOLID principles changed our design from the first component list to the final architecture, which problems we met, and how the result became more modular, maintainable, and extensible.

\Needspace{9\baselineskip}
\subsection*{8.1. How SOLID Improved the Design}

\phantomsection
\addcontentsline{toc}{subsection}{8.1. How SOLID Improved the Design}


Each principle changed a concrete part of the design:

\begin{itemize}[leftmargin=36pt, itemsep=2pt, topsep=4pt, parsep=2pt]
  \item \textbf{SRP split overloaded components.} In the first component list, Access Decision both checked access and drove the gate, and one attendance component both reviewed cases and changed records. Splitting them by reason to change gave Gate Transaction Processor, Access Decision, Attendance Review, and Attendance Consolidation (Table 4.2). Each one now has a clear owner and a clear test target.
  \item \textbf{DIP and OCP let us design around unknown APIs.} The controller and HR/Payroll APIs are not fully specified (I-01, I-02, I-05). Because Gate Core and Payroll Handover define ITurnstileControl and IPayrollSubmission themselves, we could design the gate and handover rules against clear meanings and leave the vendor details to the adapters.
  \item \textbf{LSP made us define what ``confirmed'' means.} Writing down the behavior every controller adapter must keep showed that a command receipt, an opening, a passage, and a closing are different facts. This became the core of ADR-02: attendance follows passage, not the opening signal.
  \item \textbf{ISP turned interface design into permission rules.} Splitting GateCoreApi by consumer showed that the kiosk only needs to submit and query scans. Together with the permission checks on the server, this narrow interface keeps the kiosk away from recordIntervention and emergency operations (Section 5.5).
\end{itemize}

\Needspace{9\baselineskip}
\subsection*{8.2. Challenges}

\phantomsection
\addcontentsline{toc}{subsection}{8.2. Challenges}


Applying the principles was not always straightforward:

\begin{itemize}[leftmargin=36pt, itemsep=2pt, topsep=4pt, parsep=2pt]
  \item \textbf{Too many interfaces.} SOLID can lead to an interface for every class. We added an abstraction only where it supports substitution, testing, or isolation of a real dependency, and kept other calls inside a module direct.
  \item \textbf{Modules versus classes.} Most examples of SOLID are about classes, but our main decisions are about modules in three applications. OCP at this level still needs a new release when an adapter changes; the gain is that the change stays in one place.
  \item \textbf{Contracts built on assumptions.} The ports had to be designed before the real APIs were known. The assumed controller and HR/Payroll signals are recorded as I-01, I-02, and I-05 and must be checked with the vendors, while the other assumptions are business rules or team policies that TMC has to approve.
  \item \textbf{Keeping repeated rules consistent.} Review exposed differences between summaries of the same rule. We use the detailed flows and operating policies as the reference, and shorten repeated explanations elsewhere instead of maintaining several full versions.
  \item \textbf{Boundaries versus the 2-second limit.} More boundaries could add delay on the gate path. Keeping the boundaries in-process inside each application kept the separation without extra network calls.
\end{itemize}

\Needspace{9\baselineskip}
\subsection*{8.3. Effect on Modularity, Maintainability, and Extensibility}

\phantomsection
\addcontentsline{toc}{subsection}{8.3. Effect on Modularity, Maintainability, and Extensibility}


The result is easier to change than our first draft:

\begin{itemize}[leftmargin=36pt, itemsep=2pt, topsep=4pt, parsep=2pt]
  \item \textbf{Modularity.} Thirteen modules have distinct responsibilities and communicate through public contracts. Data-owning modules control their own data; modules such as Gate Intervention and Reporting use those contracts without owning business records. The dependencies have no cycles (Table 4.4).
  \item \textbf{Maintainability.} A change usually starts in one module: a new controller in Turnstile Adapter, a new HR/Payroll format in HR/Payroll Adapter, and a new aggregation rule in Work Schedule Policy (Table 7.1).
  \item \textbf{Extensibility.} A new factory is mostly deployment and configuration, because each factory runs the same application with its own data and token scope.
  \item \textbf{What the code has shown so far.} The Task 2 prototype (Chapter 9) gave us a first check of these claims. The import tests show that the business layer of Gate Core does not depend on storage, HTTP, or adapters. Because the ports were fixed first, Packages 1, 3, and 6 were written and tested on fakes before the SQLite repositories and the turnstile adapter existed, which is the DIP benefit we expected. \tofill{Add the LSP result of the shared turnstile contract test (Package 5) and the SQLite result (Package 4).} Maintainability over a longer period, and substitution with a real controller, are still design claims.
\end{itemize}
